% Part: lambda-calculus
% Chapter: introduction
% Section: unique-readability

\documentclass[../../../include/open-logic-section]{subfiles}

\begin{document}

\olfileid[hi]{lam}{syn}{unq}
\olsection{अद्वितीय पठनीयता}

हम पूछ सकते हैं कि क्या प्रत्येक पद को बनाने की विधि अद्वितीय है---और वह
है। प्रत्येक लैम्ब्डा पद के निर्माण और निर्वचन की केवल एक विधि है। आगे की
चर्चा में \emph{निर्माण} से हमारा आशय \olref[trm]{defn:term} के निर्माण
नियमों का एक या अनेक बार उपयोग करके कोई पद बनाने की प्रक्रिया से है।

\begin{lem}\ollabel{lem:term-start}
कोई पद या तो किसी चर से आरंभ होता है या किसी कोष्ठक से।
\end{lem}

\begin{proof}
कोई वस्तु पद तभी मानी जाती है जब उसका निर्माण \olref[trm]{defn:term} के
अनुसार हुआ हो। यदि वह \olref[trm]{defn:term-var} का परिणाम है, तो उसे चर
होना चाहिए। यदि वह \olref[trm]{defn:term-abs} अथवा
\olref[trm]{defn:term-app} का परिणाम है, तो वह कोष्ठक से आरंभ होती है।
\end{proof}

\begin{lem}\ollabel{lem:app-start}
किसी अनुप्रयोग का परिणाम या तो दो कोष्ठकों से आरंभ होता है, या एक कोष्ठक
और एक चर से।
\end{lem}

\begin{proof}
यदि $M$ किसी अनुप्रयोग का परिणाम है, तो वह $(PQ)$ रूप का है और इसलिए
कोष्ठक से आरंभ होता है। चूँकि $P$ कोई पद है, इसलिए
\olref{lem:term-start} से वह या तो कोष्ठक से अथवा चर से आरंभ होता है।
\end{proof}

\begin{lem}\ollabel{lem:initial}
किसी पद का कोई उचित आरंभिक भाग स्वयं पद नहीं होता।
\end{lem}

\begin{prob}
पदों की लंबाई पर आगमन द्वारा \olref[lam][syn][unq]{lem:initial} सिद्ध
करें।
\end{prob}

\begin{prop}[अद्वितीय पठनीयता] \ollabel{prop:unq}
प्रत्येक पद का निर्माण अद्वितीय है। दूसरे शब्दों में, यदि पद $M$ किसी
निर्माण द्वारा बनाया गया है, तो वही एकमात्र निर्माण है जो इस पद को बना
सकता है।
\end{prop}

\begin{proof}
  हम इसे पदों के निर्माण पर आगमन द्वारा सिद्ध करते हैं।
  \begin{enumerate}
    \item $M$, ~$x$ रूप का है, जहाँ $x$ कोई चर है। चूँकि अमूर्तनों और
      अनुप्रयोगों के परिणाम सदा कोष्ठकों से आरंभ होते हैं, इसलिए $M$ के
      निर्माण में उनका उपयोग नहीं हुआ हो सकता। अतः $M$ का निर्माण
      \olref[trm]{defn:term}\olref[trm]{defn:term-var} का एक ही चरण होना
      चाहिए।
    \item $M$, ~$(\lambd[x][N])$ रूप का है, जहाँ $x$ कोई चर और $N$ कोई
      पद है। उसका निर्माण
      \olref[trm]{defn:term}\olref[trm]{defn:term-var} के अनुसार नहीं हुआ
      हो सकता, क्योंकि वह अकेला चर नहीं है। \olref{lem:app-start} से वह
      किसी अनुप्रयोग का परिणाम भी नहीं है। अतः $M$ केवल~$N$ पर अमूर्तन का
      परिणाम हो सकता है। आगमन-परिकल्पना से हम जानते हैं कि $N$ का निर्माण
      स्वयं अद्वितीय है।
    \item $M$, ~$(PQ)$ रूप का है, जहाँ $P$ और $Q$ पद हैं। चूँकि वह कोष्ठक
      से आरंभ होता है, इसलिए उसका निर्माण
      \olref[trm]{defn:term}\olref[trm]{defn:term-var} द्वारा भी नहीं हो
      सकता। \olref{lem:term-start} से $P$, ~$\lambd$ से आरंभ नहीं हो
      सकता, अतः $(PQ)$ किसी अमूर्तन का परिणाम नहीं हो सकता। अब मान लें कि
      अनुप्रयोग द्वारा $M$ बनाने की कोई दूसरी विधि भी होती, जैसे वह
      $(P'Q')$ रूप का भी होता। तब $P$, ~$P'$ का उचित आरंभिक खंड होता (या
      उलटा), जो \olref{lem:initial} से असंभव है। अतः $P$ और $Q$ अद्वितीय
      रूप से निर्धारित हैं, और आगमन-परिकल्पना से हम जानते हैं कि उनके
      निर्माण भी अद्वितीय हैं।
  \end{enumerate}
\end{proof}

ऊपर के साध्य का अधिक पठनीय पुनर्कथन इस प्रकार है:
\begin{prop}
  पद $M$ केवल निम्न रूपों में से किसी एक का हो सकता है:
  \begin{enumerate}
    \item $x$, जहाँ $x$ ऐसा चर है जो $M$ से अद्वितीय रूप से निर्धारित है।
    \item $(\lambd[x][N])$, जहाँ $x$ कोई चर और $N$ कोई दूसरा पद है; दोनों
      $M$ से अद्वितीय रूप से निर्धारित हैं।
    \item $(PQ)$, जहाँ $P$ और $Q$ ऐसे दो पद हैं जो $M$ से अद्वितीय रूप से
      निर्धारित हैं।
  \end{enumerate}
\end{prop}

\end{document}
